\chapter{The Kalman filter}\label{ch:kf}

This note documents \code{KalmanFilter}, the \code{nabia} forecaster for linear
Gaussian state-space models with a scalar observation. We first prove the Gaussian
conditioning lemma and derive the filter from it. We then prove that the Kalman gain
minimizes the error covariance among all linear updates, in the Loewner order, and
show why the Joseph form of the covariance update is used. We derive the
likelihood by the prediction error decomposition, the $h$-step forecast
distribution, the treatment of missing values and the Rauch--Tung--Striebel
smoother. For the two built-in models, local level and local linear trend, we
derive the steady state, the link with exponential smoothing and the growth of
the forecast variance with the horizon. The last section maps the algorithm to
the code and discusses numerical choices, complexity and limitations.




%==============================================================================
\section{Model}

A linear Gaussian state-space model with $n$ states and a scalar observation is
\begin{equation}\label{kf:eq:model}
\begin{aligned}
  x_t &= F x_{t-1} + w_t, & w_t &\sim \N(0, Q),\\
  y_t &= H x_t + v_t,     & v_t &\sim \N(0, R),
\end{aligned}
\qquad t = 0,\dots,T-1,
\end{equation}
with transition $F\in\R^{n\times n}$, observation row $H\in\R^{1\times n}$,
process noise covariance $Q\succeq0$ and observation noise variance $R>0$. The
prior on the first state is $x_0\sim\N(\bar x_0, P_0)$. The initial state, the
$w_t$ and the $v_t$ are mutually independent and independent over time. The state
$x_t$ is hidden; only $y_t$ is observed, and some $y_t$ may be missing.

\paragraph{Notation.} For $s\le t$, write $y_{0:s}$ for the observations up
to time $s$ and
\[
  x_{t\mid s} = \E[x_t\mid y_{0:s}],\qquad P_{t\mid s} = \Cov(x_t\mid y_{0:s}).
\]
With this notation, $x_{0\mid -1}=\bar x_0$ and $P_{0\mid -1}=P_0$. The \emph{filter} computes
$(x_{t\mid t},P_{t\mid t})$, the \emph{one-step predictor} $(x_{t\mid t-1},P_{t\mid t-1})$,
the \emph{forecast} $(x_{T-1+h\mid T-1},P_{T-1+h\mid T-1})$ and the \emph{smoother}
$(x_{t\mid T-1},P_{t\mid T-1})$.

%==============================================================================
\section{The filter}

\subsection{Gaussian conditioning}

\begin{lemma}[Conditioning a Gaussian vector]\label{kf:lem:conditioning}
Let $(a,b)$ be jointly Gaussian with means $\mu_a,\mu_b$, covariances
$\Sigma_{aa},\Sigma_{bb}$, cross-covariance $\Sigma_{ab}=\Sigma_{ba}^\T$, and
$\Sigma_{bb}\succ0$. Then
\[
  a\mid b \;\sim\; \N\!\big(\mu_a + \Sigma_{ab}\Sigma_{bb}^{-1}(b-\mu_b),\;
                         \Sigma_{aa}-\Sigma_{ab}\Sigma_{bb}^{-1}\Sigma_{ba}\big).
\]
\end{lemma}
\begin{proof}
Let $G=\Sigma_{ab}\Sigma_{bb}^{-1}$ and $c=a-Gb$. Then
$\Cov(c,b)=\Sigma_{ab}-G\Sigma_{bb}=0$. Since $(c,b)$ is a linear function of
the Gaussian vector $(a,b)$, it is jointly Gaussian, and zero covariance implies
that $c$ and $b$ are independent. Writing $a=c+Gb$, conditionally on $b$ the
vector $a$ is $c$ (whose distribution is unchanged by conditioning) plus the
constant $Gb$. Hence $a\mid b$ is Gaussian with mean
$\E c+Gb=\mu_a-G\mu_b+Gb$ and covariance
$\Cov(c)=\Sigma_{aa}-G\Sigma_{ba}-\Sigma_{ab}G^\T+G\Sigma_{bb}G^\T=\Sigma_{aa}-\Sigma_{ab}\Sigma_{bb}^{-1}\Sigma_{ba}$.
\end{proof}

\subsection{Prediction and correction}

\begin{theorem}[Kalman filter]\label{kf:thm:kalman}
For model \eqref{kf:eq:model}, every distribution $x_t\mid y_{0:s}$ is Gaussian, and the
moments obey, for $t\ge 0$:
\begin{align}
  \text{predict ($t\ge1$):}\quad & x_{t\mid t-1} = F x_{t-1\mid t-1}, &
     P_{t\mid t-1} &= F P_{t-1\mid t-1} F^\T + Q, \label{kf:eq:predict}\\
  \text{innovation:}\quad & e_t = y_t - H x_{t\mid t-1}, &
     S_t &= H P_{t\mid t-1} H^\T + R, \label{kf:eq:innovation}\\
  \text{gain:}\quad & K_t = P_{t\mid t-1} H^\T / S_t, \label{kf:eq:gain}\\
  \text{correct:}\quad & x_{t\mid t} = x_{t\mid t-1} + K_t e_t, &
     P_{t\mid t} &= (I-K_tH)\,P_{t\mid t-1}. \label{kf:eq:correct}
\end{align}
Moreover $y_t\mid y_{0:t-1}\sim\N(Hx_{t\mid t-1},S_t)$.
\end{theorem}
\begin{proof}
By induction on $t$. Assume $x_{t-1}\mid y_{0:t-1}\sim\N(x_{t-1\mid t-1},P_{t-1\mid t-1})$.
For $t=0$, the prior plays this role and there is no prediction step.

\emph{Prediction.} The noise $w_t$ is independent of $(x_{t-1},y_{0:t-1})$. Given
$y_{0:t-1}$, the state $x_t=Fx_{t-1}+w_t$ is an affine function of two
independent Gaussian vectors. So $x_t\mid y_{0:t-1}$ is Gaussian with mean
$Fx_{t-1\mid t-1}$ and covariance $FP_{t-1\mid t-1}F^\T+Q$.

\emph{Correction.} The noise $v_t$ is independent of $(x_t,y_{0:t-1})$, so given
$y_{0:t-1}$ the pair $(x_t,y_t)$ is jointly Gaussian with
\[
  \E[y_t\mid y_{0:t-1}] = Hx_{t\mid t-1},\qquad
  \Var(y_t\mid y_{0:t-1}) = HP_{t\mid t-1}H^\T+R = S_t,
\]
\[
  \Cov(x_t,y_t\mid y_{0:t-1}) = P_{t\mid t-1}H^\T .
\]
This already proves the claim on $y_t\mid y_{0:t-1}$. Conditioning
additionally on $y_t$ means applying Lemma~\ref{kf:lem:conditioning} to this
conditional joint distribution, with $a=x_t$, $b=y_t$ and $\Sigma_{bb}=S_t\ge R>0$.
This gives the mean $x_{t\mid t-1}+P_{t\mid t-1}H^\T S_t^{-1}e_t$ and covariance
$P_{t\mid t-1}-P_{t\mid t-1}H^\T S_t^{-1}HP_{t\mid t-1}=(I-K_tH)P_{t\mid t-1}$.
\end{proof}

With a scalar observation, $S_t$ is a positive number. The gain is a vector,
and no matrix is ever inverted.

\subsection{Optimality of the gain and the Joseph form}

Theorem~\ref{kf:thm:kalman} uses Gaussianity. Without it, the filter is still the
best \emph{linear} update.

\begin{proposition}[Joseph form, valid for any gain]\label{kf:prop:joseph}
Let $\hat x^-$ be an estimate of $x$ with error covariance
$P^-=\Cov(x-\hat x^-)$, and observe $y=Hx+v$ with $v$ zero-mean, variance $R$,
and uncorrelated with $x-\hat x^-$. For any gain $G\in\R^{n}$, the update
$\hat x^+=\hat x^-+G(y-H\hat x^-)$ has error covariance
\begin{equation}\label{kf:eq:joseph}
  P^+(G) = (I-GH)\,P^-\,(I-GH)^\T + G R G^\T .
\end{equation}
\end{proposition}
\begin{proof}
The new error is $x-\hat x^+ = (x-\hat x^-) - G(Hx+v-H\hat x^-) = (I-GH)(x-\hat x^-) - Gv$.
The two terms are uncorrelated, so their covariances add, giving \eqref{kf:eq:joseph}.
\end{proof}

\begin{theorem}[The Kalman gain is Loewner-optimal]\label{kf:thm:optimal}
With $S=HP^-H^\T+R$ and $K=P^-H^\T/S$, for every $G$,
\[
  P^+(G) = (I-KH)P^- + (G-K)\,S\,(G-K)^\T \;\succeq\; P^+(K) = (I-KH)P^- .
\]
In particular $K$ minimizes $\tr P^+(G)$ (the mean squared error) and every
other weighted mean squared error $\E\norm{x-\hat x^+}_W^2$, $W\succeq0$.
\end{theorem}
\begin{proof}
Expanding \eqref{kf:eq:joseph},
$P^+(G)=P^- - GHP^- - P^-H^\T G^\T + G S G^\T$. Using $P^-H^\T=KS$ and
completing the square,
\[
  P^+(G) = P^- - KSK^\T + (G-K)S(G-K)^\T .
\]
Since $KSK^\T = KHP^-$, the first two terms equal $(I-KH)P^-$. The last term is positive
semi-definite because $S>0$, and it vanishes iff $G=K$. For the weighted error,
$\E\norm{e}_W^2=\tr(WP^+)$ and $\tr(W(G-K)S(G-K)^\T)\ge0$.
\end{proof}

\paragraph{Why the implementation uses the Joseph form.}
The short form $(I-KH)P^-$ equals \eqref{kf:eq:joseph} only for the exact gain.
In floating point the computed $K$ is slightly off, and the short form then
loses symmetry and can even lose positive semi-definiteness: it subtracts two
nearly equal matrices. Equation \eqref{kf:eq:joseph} is a sum of two positive semi-definite
terms, so it stays symmetric positive semi-definite for any $K$. By
Theorem~\ref{kf:thm:optimal}, an error $\delta$ in the gain changes it only by
$\delta S\delta^\T=O(\norm{\delta}^2)$. \code{nabia} therefore updates the
covariance as $P\gets(I-KH)P(I-KH)^\T+R\,KK^\T$.

\subsection{Missing observations}
If $y_t$ is missing, conditioning on the available data is conditioning on
$y_{0:t-1}$ only. Then $x_{t\mid t}=x_{t\mid t-1}$ and $P_{t\mid t}=P_{t\mid t-1}$: the
correction step is skipped. Uncertainty grows through consecutive prediction steps
and shrinks again at the next observation. \code{nabia} treats any non-finite
$y_t$ (NaN) as missing.

%==============================================================================
\section{Likelihood}

\begin{proposition}[Prediction error decomposition]\label{kf:prop:likelihood}
Let $\mathcal O\subseteq\{0,\dots,T-1\}$ be the set of observed times. The log
density of the observed values is
\[
  \log p\big((y_t)_{t\in\mathcal O}\big) = -\frac12\sum_{t\in\mathcal O}\Big(\log(2\pi) + \log S_t + \frac{e_t^2}{S_t}\Big).
\]
\end{proposition}
\begin{proof}
By the chain rule of probability, the joint density of the observed values is the product of
the conditional densities $p(y_t\mid \text{observed }y_s,\ s<t)$ over
$t\in\mathcal O$. By Theorem~\ref{kf:thm:kalman}, applied with the missing values
skipped, each factor is $\N(y_t;\,Hx_{t\mid t-1},S_t)$. Taking logs gives the result.
\end{proof}

The likelihood therefore comes for free with the filter, in $O(1)$ extra work
per step. It is the objective for estimating unknown entries of $F,H,Q,R$ by
maximum likelihood.

\begin{remark}[Diffuse prior]
With a large $P_0$ (the \code{nabia} default is $10^7$), the first term contains
$\log S_0\approx\log P_0$. This term reflects the vague prior, not the
data. When comparing likelihoods across parameter values with the same $P_0$,
the term is nearly constant and harmless. To compare models with
different state dimensions, drop the first $n$ terms (the diffuse
likelihood).
\end{remark}

%==============================================================================
\section{Forecasting}

\begin{proposition}[$h$-step forecast]\label{kf:prop:forecast}
For $h\ge1$, $y_{T-1+h}\mid y_{0:T-1}\sim\N(\mu_h,v_h)$ with
\begin{align*}
  x_{T-1+h\mid T-1} &= F^h x_{T-1\mid T-1}, & P_{T-1+h\mid T-1} &= F P_{T-2+h\mid T-1} F^\T + Q,\\
  \mu_h &= Hx_{T-1+h\mid T-1}, & v_h &= HP_{T-1+h\mid T-1}H^\T+R .
\end{align*}
\end{proposition}
\begin{proof}
Future observations are missing data. Apply the prediction step
\eqref{kf:eq:predict} $h$ times without corrections, then take the observation
moments as in the proof of Theorem~\ref{kf:thm:kalman}.
\end{proof}

%==============================================================================
\section{Smoothing}

The filter uses only the past. For analysis (for example, extracting the trend
of a historical series), the smoother uses the whole sample.

\begin{theorem}[Rauch--Tung--Striebel smoother~\cite{rts1965}]\label{kf:thm:rts}
Let $x_{T-1\mid T-1}, P_{T-1\mid T-1}$ be the last filtered moments. For
$t=T-2,\dots,0$, with $C_t = P_{t\mid t}F^\T P_{t+1\mid t}^{-1}$,
\[
  x_{t\mid T-1} = x_{t\mid t} + C_t\,(x_{t+1\mid T-1}-x_{t+1\mid t}),\qquad
  P_{t\mid T-1} = P_{t\mid t} + C_t\,(P_{t+1\mid T-1}-P_{t+1\mid t})\,C_t^\T .
\]
\end{theorem}
\begin{proof}
\emph{Step 1 (Markov property).} Given $x_{t+1}$, the state $x_t$ is
independent of $y_{t+1:T-1}$. Those observations are functions of $x_{t+1}$ and
of the noises $w_{t+2:T-1},v_{t+1:T-1}$, which are independent of $x_t$ and
$y_{0:t}$. Hence $p(x_t\mid x_{t+1},y_{0:T-1})=p(x_t\mid x_{t+1},y_{0:t})$.

\emph{Step 2.} Given $y_{0:t}$, the pair $(x_t,x_{t+1})$ is jointly Gaussian with
$\Cov(x_t,x_{t+1})=P_{t\mid t}F^\T$ and $\Cov(x_{t+1})=P_{t+1\mid t}$.
Lemma~\ref{kf:lem:conditioning} gives
$x_t\mid x_{t+1},y_{0:t}\sim\N\big(x_{t\mid t}+C_t(x_{t+1}-x_{t+1\mid t}),\ P_{t\mid t}-C_tP_{t+1\mid t}C_t^\T\big)$.

\emph{Step 3.} Average over $x_{t+1}\mid y_{0:T-1}\sim\N(x_{t+1\mid T-1},P_{t+1\mid T-1})$.
The law of total expectation gives the mean. The law of total variance gives
$P_{t\mid t}-C_tP_{t+1\mid t}C_t^\T + C_tP_{t+1\mid T-1}C_t^\T$, which is the
stated covariance.
\end{proof}

\code{KalmanFilter::smoothed\_states()} returns the smoothed means. It computes
$C_t$ as $\big(P_{t+1\mid t}^{-1}FP_{t\mid t}\big)^\T$ with an $LDL^\T$ solve,
using the symmetry of both covariances, rather than forming an inverse.

%==============================================================================
\section{Built-in models}

\subsection{Local level}
$F=H=1$, $Q=q$, $R=r$: a random walk observed with noise. The predicted variance
follows the scalar Riccati recursion
$P_{t+1\mid t} = P_{t\mid t-1}r/(P_{t\mid t-1}+r) + q$.

\begin{proposition}[Steady state]\label{kf:prop:ll-steady}
The predicted variance converges to $P=\tfrac12\big(q+\sqrt{q^2+4qr}\big)$ and
the gain to $K=P/(P+r)$. The filter then becomes simple exponential smoothing with
$\alpha=K$, and $q/r=K^2/(1-K)$.
\end{proposition}
\begin{proof}
A fixed point solves $P^2-qP-qr=0$, whose positive root is the stated $P$. The map
$g(P)=Pr/(P+r)+q$ is increasing with $g'(P)=r^2/(P+r)^2<1$ for $P>0$, and it is
concave. Hence the iterates are monotone and bounded, and they converge to the
unique positive fixed point from any $P_0\ge0$. With $F=H=1$, \eqref{kf:eq:correct}
reads $x_{t\mid t}=x_{t-1\mid t-1}+K_t(y_t-x_{t-1\mid t-1})$, the SES recursion. The
relation $q/r=K^2/(1-K)$ follows from substituting $P=Kr/(1-K)$ into
$P^2=q(P+r)$.
\end{proof}

\paragraph{Diffuse start.}
With $P_0\gg r$, the first gain is $K_0=P_0/(P_0+r)\approx1$, so
$x_{0\mid0}\approx y_0$ and $P_{0\mid0}=P_0r/(P_0+r)\approx r$. A large initial
variance thus amounts to starting from the first observation, as SES does.

\subsection{Local linear trend}
With state $x_t=(\ell_t,b_t)$ (level and slope),
\[
  F=\begin{pmatrix}1&1\\0&1\end{pmatrix},\quad H=\begin{pmatrix}1&0\end{pmatrix},\quad
  Q=\diag(q_\ell,q_b).
\]

\begin{proposition}[Trend extrapolation and variance growth]\label{kf:prop:llt}
The point forecast is the straight line $\mu_h=\ell_{T-1\mid T-1}+h\,b_{T-1\mid T-1}$. The
part of the forecast variance due to future slope shocks equals
$q_b\sum_{m=1}^{h-1}m^2=q_b\,(h-1)h(2h-1)/6\sim q_bh^3/3$.
\end{proposition}
\begin{proof}
By induction, $F^h=\begin{pmatrix}1&h\\0&1\end{pmatrix}$, since
$F^{h+1}=F^hF=\begin{pmatrix}1&h+1\\0&1\end{pmatrix}$. With
Proposition~\ref{kf:prop:forecast}, this gives the mean. For the variance, unroll the model from $T-1$:
$b_{T-1+k}=b_{T-1}+\sum_{i=1}^{k}w^b_{T-1+i}$, and
\[
  \ell_{T-1+h} = \ell_{T-1} + \sum_{k=0}^{h-1} b_{T-1+k} + \sum_{j=1}^{h} w^\ell_{T-1+j}
               = \ell_{T-1} + h\,b_{T-1} + \sum_{i=1}^{h-1}(h-i)\,w^b_{T-1+i} + \sum_{j=1}^{h}w^\ell_{T-1+j}.
\]
The future shocks are independent of the current state. The slope shocks contribute
$q_b\sum_{i=1}^{h-1}(h-i)^2=q_b\sum_{m=1}^{h-1}m^2$.
\end{proof}

So a trend model extrapolates the current slope. Its uncertainty grows like $h^3$,
much faster than the linear growth of the local level model, because any
error in the slope accumulates over the horizon. A small $q_b$ gives a nearly
deterministic trend, and $q_b=0$ gives a fixed slope estimated from the data.

%==============================================================================
\section{Implementation}

\subsection{Interface}
\begin{lstlisting}
auto p = std::make_unique<KalmanFilterParams>(
    KalmanFilterParams::local_linear_trend(/*q_level*/ 1e-4, /*q_slope*/ 1e-8, /*R*/ 0.09));
KalmanFilter kf(std::move(p));
kf.fit(y);                                   // NaN = missing
Forecast fc = kf.forecast(10);               // mean and variance (Prop. 4.1)
const Eigen::MatrixXd& x = kf.filtered_states();   // n x T, column t = x_{t|t}
Eigen::MatrixXd xs = kf.smoothed_states();         // n x T, column t = x_{t|T-1}
double ll = kf.log_likelihood();                   // Prop. 3.1
\end{lstlisting}
To define any other model, fill in the parameter fields:
\begin{center}
\begin{tabular}{@{}llll@{}}
\code{transition} & $F$ & \code{observation\_noise} & $R$ \\
\code{observation} & $H$ & \code{initial\_state} & $\bar x_0$ \\
\code{process\_noise} & $Q$ & \code{initial\_covariance} & $P_0$ \\
\end{tabular}
\end{center}

\subsection{Algorithm}
\begin{algorithm}[ht]
\caption{\code{KalmanFilter::fit}}\label{kf:alg:fit}
\begin{algorithmic}[1]
\State $x\gets\bar x_0$;\; $P\gets P_0$;\; $\ell\gets0$
\For{$t=0,\dots,T-1$}
  \If{$t>0$} $x\gets Fx$;\; $P\gets FPF^\T+Q$ \EndIf \Comment{predict}
  \State store $x_{t\mid t-1}=x$, $P_{t\mid t-1}=P$
  \State $\hat y_t\gets Hx$;\; $u\gets PH^\T$;\; $S_t\gets Hu+R$ \Comment{$u$ reused by the gain}
  \If{$y_t$ is finite}
    \State $e_t\gets y_t-\hat y_t$;\; $K\gets u/S_t$;\; $x\gets x+Ke_t$
    \State $A\gets I-KH$;\; $P\gets APA^\T+R\,KK^\T$ \Comment{Joseph form}
    \State $\ell\gets\ell-\tfrac12(\log2\pi+\log S_t+e_t^2/S_t)$
  \Else\ $e_t\gets$ NaN \Comment{missing: no correction}
  \EndIf
  \State store $x_{t\mid t}=x$, $P_{t\mid t}=P$
\EndFor
\end{algorithmic}
\end{algorithm}

\paragraph{Complexity.} Each step costs $O(n^3)$ for the products with $F$ and
$O(n^2)$ for the rank-one correction, so fitting is $O(Tn^3)$. The filter stores the
predicted and filtered moments ($O(Tn^2)$ memory) because the smoother needs
them. Forecasting costs $O(hn^3)$, and smoothing $O(Tn^3)$. For the built-in models
$n\le2$, and the cost is linear in $T$.

\paragraph{Validation.} The constructor checks that $F,H,Q,\bar x_0,P_0$ have consistent
dimensions, that $Q$ and $P_0$ are square and symmetric (to a relative tolerance of $10^{-10}$),
and that $R>0$ is finite. It throws \code{std::invalid\_argument} otherwise.
\code{fit} rejects an empty series. Accessors called before \code{fit} throw
\code{std::logic\_error}.

\paragraph{Correspondence with the code.}
\begin{center}
\begin{tabular}{@{}ll@{}}
\toprule
Mathematics & \code{KalmanFilter} \\
\midrule
$x_{t\mid t},\ P_{t\mid t}$ & \code{filtered\_states().col(t)}, \code{filtered\_covariances()[t]} \\
$Hx_{t\mid t-1}$ & \code{fitted\_values()(t)} \\
$e_t,\ S_t$ & \code{innovations()(t)}, \code{innovation\_variances()(t)} \\
$\log p(y)$ & \code{log\_likelihood()} \\
$x_{t\mid T-1}$ & \code{smoothed\_states().col(t)} \\
$\mu_h,\ v_h$ & \code{forecast(H).mean(h-1)}, \code{.variance(h-1)} \\
\bottomrule
\end{tabular}
\end{center}

\paragraph{Tests.} \code{tests/test\_forecasters.cpp} checks:
\begin{itemize}
  \item Proposition~\ref{kf:prop:ll-steady}: the filter matches SES with the steady-state gain to $10^{-9}$.
  \item Proposition~\ref{kf:prop:likelihood}: the log-likelihood equals the log
        density of $y\sim\N(\bar x_0\mathbf1,\Sigma)$ with
        $\Sigma_{ij}=P_0+q\min(i,j)+r\,\delta_{ij}$, computed directly, to $10^{-10}$.
  \item Proposition~\ref{kf:prop:llt}: the slope state recovers a true slope of 0.5,
        and forecasts lie on the extrapolated line.
  \item Proposition~\ref{kf:prop:forecast}: for the local level, $v_h=P_{T-1\mid T-1}+hq+r$.
  \item Missing values: the variance grows at missing steps and shrinks again at the next observation.
  \item Smoother: with $q=0$, every smoothed state equals the final filtered estimate.
\end{itemize}

\subsection{Limitations and extensions}
\begin{itemize}
  \item Observations are scalar and the system matrices are time-invariant. Vector
        observations would replace the division by $S_t$ with a Cholesky solve.
  \item $Q$, $R$ (and any unknowns in $F,H$) are inputs. A maximum-likelihood fit
        can maximize \code{log\_likelihood()} over log-variances, for example with
        \code{NelderMead} or, using analytic score recursions, with \code{LBFGS}.
  \item The smoother returns means only. Theorem~\ref{kf:thm:rts} also gives the
        smoothed covariances if needed.
\end{itemize}

%==============================================================================
